\documentclass[a4paper]{article}
\usepackage{color}
\usepackage{graphicx}
\usepackage[left=1.5cm,top=1cm,right=1.5cm,bottom=1cm]{geometry}
\graphicspath{ {Results/} }
\usepackage[paper=portrait,pagesize]{typearea}
\usepackage{multirow}
\usepackage{dblfloatfix}
\usepackage{caption}
\usepackage{parselines} 
\usepackage{subcaption}
\usepackage{multicol}
\usepackage{xfrac} 
\usepackage{nicefrac}
\usepackage{subfiles}
\usepackage{gensymb}
\usepackage{wrapfig,float}
\setlength{\hoffset}{-0.5in}\hoffset-0.5in
\setlength{\textwidth}{15cm}
\topmargin = 20pt
\voffset = -20pt
\addtolength{\textheight}{2cm}
\author{Abigail Millward}
\date{8th of November 2017}
    
%--------------------------------
\usepackage[
    backend=biber,
    style=authoryear,
    url=false,
    isbn=false,
    doi=false,
    dateabbrev=false,
    language=british
    ]{biblatex}
     
\addbibresource{Code/project_pro.bib}
    
\begin{document}

\thispagestyle{empty}
\null\vskip0.2in%
\begin{center}
\LARGE{{\bf 
Classificaiton of complex stresses in plants using machine learning}}
\end{center}

\vspace{0.5cm}

\begin{center}
{\Large {\bf by}}\\
\mbox{} \\
{\Large {\bf Abigail Baines (CID: 00742373)}}

{\Large Supervisor : Oliver Windram, Imperial College London, o.windram@imperial.ac.uk}
\end{center}

\vspace{1cm}

\begin{center}
\large{\bf{Department of Life Sciences \\ Imperial College London \\
London SL5 7PY \\ United Kingdom}}
\end{center}


\vspace{1cm}

\begin{figure}[!h]
\centering
\includegraphics[scale=0.35]{imperial_crest_colour.jpg}
\end{figure}

\vspace{1cm}

\begin{center}
\large{\bf{Thesis submitted as part of the requirements for the award of the 
MRes in Computational Methods in Ecology and Evolution, Imperial College London, 2017-2018}} 
\vspace{0.5cm}
\large{\bf{

Formatted in the style of frontiers in Plant Science}}
\vspace{0.5cm}
\large{\bf{

Word count: 5455}}

\end{center}


\vspace{2cm}



\thispagestyle{empty}
\null\vskip0.2in%
\begin{center}
\LARGE{{\bf 
Declarations:}}
\end{center}

\vspace{0.5cm}
\begin{center}
{\Large Data provided by Dr Oliver Windram, Chris Adams and Abigail Baines

Data processing, cleaning and augmenting by Abigail Baines

All mathematical models developed by Abigail Baines

All interpretations of results and analyses were developed by Abigail Baines}
\end{center}
\pagebreak

\begin{abstract}
      Determining the incidence and types of plant stresses has been of interest for many years, particularly
    due to their negative impact on agricultural production. Despite this, current methods
    are very qualitative, with many relying on manual detection which can be highly variable.
    With this in mind, our aim for this project was to see whether we could use machine learning, in particular Random Forests (RFs) and Convoltional Neural Networks (CNNs),
    to classify plant stresses from RGB and colour-pass filtered leaf images. Two plant leaf datasets were collected; an infection duration study conducted to assess the impact of \textit{Botrytis cinerea} on \textit{Arabadopsis thaliana} over a 72 hour period, and a complex stress study conducted on tomato leaves. Results showed support for four out of the five proposed hyptheses; hypothesis four (normalised images will improve model performance) was not supported by the data collected. We found strong support that machine learning  can correctly identify both infection duration and complex/multi stresses in an image; both Rfs and CNNs produced training accuracies of \textgreater 85\%. Finally, it was noted that the use of colour-pass filters improve model performance. Overall the results obtained from this study will be greatly benefical to both the botantical and horticultural communities, and to the future of crop production. 
    
    \vspace{5mm}
    Keywords: Digital Pathology, Random Forests, Python, Disease detection,
    Multi-spectral features, Machine Learning.
  \end{abstract}
         
\section{Introduction}

\subfile{Code/Sections/Introduction}

\section{Materials \& Methods}

\subfile{Code/Sections/Materials_Methods}

\section{Results}

\subfile{Code/Sections/Results}

\section{Discussion}

\subfile{Code/Sections/Discussion}

\section{Conclusions}

\subfile{Code/Sections/Conclusion}

\printbibliography

\pagebreak  
\KOMAoptions{paper=landscape,pagesize}
\recalctypearea    
\section{Appendix}
\subfile{Code/Sections/Appendix}


\end{document}